% ============================================================
%  ApJ Full Article — Surface Density as Organizing Principle of LRDs
%  Author: Tan Xin (Independent Researcher, IAIP)
%  Date: 2026-05-12
%  Engine: pdfLaTeX (aastex701)
% ============================================================
\documentclass[apj,twocolumn,longbibliography,linenumbers]{aastex701}

\usepackage{amsmath}
\usepackage{amssymb}
\usepackage{graphicx}
\usepackage{xcolor}
\usepackage{placeins}
\usepackage{booktabs}
\usepackage{url}
\hypersetup{colorlinks=true,linkcolor=blue,urlcolor=blue,citecolor=blue}

\graphicspath{{figures/}}

%% ============================================================
%%  METADATA
%% ============================================================

\shorttitle{Surface Density Organizes LRD Spectral Properties}

\begin{document}

\title{Surface Density as the Organizing Principle of Little Red Dot
Spectral Properties}

\author{Tan Xin}
\affil{Independent Researcher, IAIP}
\email{summicron352@gmail.com}
\correspondingauthor{summicron352@gmail.com}

\keywords{galaxies: high-redshift (740) --- galaxies: evolution (594) --- galaxies: photometry (611) --- gravitation (662) --- cosmology: observations (346)}

%% ============================================================
%%  ABSTRACT
%% ============================================================

\begin{abstract}
JWST has discovered hundreds of compact, red sources at $z \sim 4$--$9$, collectively known as Little Red Dots (LRDs), whose extreme compactness ($r_{\rm eff} < 300$ pc), anomalously red colors ($F_{444W}/F_{150W}$), broad emission lines (FWHM $\sim 1000$--$3000$ km s$^{-1}$), and mid-infrared silence defy standard galaxy evolution paradigms. Current explanations invoke at least three independent mechanisms (dust attenuation, AGN activity, and starburst/merger-driven compaction) that lack mutual physical connection, and no framework yet explains the internal ordering: why, among LRDs at the same redshift and luminosity, do more compact sources systematically appear redder and exhibit broader lines?

I present the first comprehensive surface density analysis of 260 photometric LRDs from Kokorev et al.\ (2024). The color--$\Sigma$ partial correlation is $\rho_p = +0.341$ ($5.6\sigma$) after controlling for redshift and bolometric luminosity---more compact LRDs are systematically redder at fixed $z$ and $L_{\rm bol}$. An exhaustive control-variable analysis of 1,626 source pairs from a combined 296-object catalog reveals that 34.6\% of matched pairs show reversed color ordering (lower-$z$, higher-$\Sigma$ sources appearing redder than higher-$z$, lower-$\Sigma$ counterparts; $p < 10^{-30}$), decisively rejecting pure Balmer break dominance. A dose--response relationship ($\rho = +0.236$, $p = 4.8 \times 10^{-22}$) between $|\Delta\Sigma|$ and $|\Delta{\rm color}|$ among reversed pairs establishes $\Sigma$ as an independent causal driver. Nine independent robustness tests (R1--R9), including Monte Carlo error propagation, per-field validation across six JWST surveys, and jackknife resampling, unanimously preserve the signal.

Independent spectroscopic verification with 36 RUBIES LRDs confirms the same physical picture: FWHM--$\Sigma$ correlation $\rho = +0.632$ ($4.1\sigma$), cross-validated by Euclid (Tumborang et al.\ 2026) and directionally confirmed in JADES DR5 despite ${\sim}50\%$ sample contamination ($\rho_p = +0.136$, $5.2\sigma$).

I propose density-dependent effective gravity, $G_{\rm eff}(\Sigma) = G_N[1 + \varepsilon_g(\Sigma/\Sigma_0)^\beta]$, as a mechanism-agnostic phenomenological framework that unifies the four LRD anomalies through a single variable $\Sigma$: deeper potential wells enhance gravitational redshift (reddening) and virial velocities (line broadening) without requiring dust thermal re-emission (MIR silence). A four-layer gravitational redshift model demonstrates that $G_{\rm eff} > G_N$ introduces a monotonic, $\Sigma$-dependent color modulation detectable through rank-order statistics, even though the absolute magnitude shift (${\sim}10^{-4}$\,mag) is small relative to the observed color excess. I present five falsifiable predictions (P1--P5) and conclude that surface density is the fundamental organizing parameter of LRD spectral properties---any complete framework must incorporate $\Sigma$ as an independent degree of freedom.
\end{abstract}

%% ============================================================
%%  §1  INTRODUCTION
%% ============================================================

\section{Introduction}

\subsection{JWST and the Little Red Dots}

Prior to JWST, the $z > 6$ universe was expected to be largely barren---a sparse population of faint, irregular, low-mass galaxies assembled from the first collapsed dark matter halos. This picture was decisively overturned within the first two JWST observing seasons. Deep NIRCam imaging from JADES, CEERS, PRIMER, and COSMOS-Web revealed hundreds of compact, red sources at $z \sim 4$--$9$, collectively dubbed Little Red Dots (LRDs; \citealt{Labbe2023}; \citealt{Kokorev2024}; \citealt{Adams2024}).

LRDs challenge standard galaxy evolutionary paradigms across multiple, seemingly independent dimensions. Their defining properties can be summarized as four interconnected anomalies:

1.~\textbf{Extreme compactness}. Effective radii $r_{\rm eff} < 300$ pc, with a substantial tail below 100 pc---one to two orders of magnitude smaller than typical star-forming galaxies at the same epoch (\citealt{Fujimoto2023}; \citealt{Kocevski2024}).

2.~\textbf{Anomalous redness}. Rest-frame optical colors, probed by the $F_{444W}/F_{150W}$ flux ratio, far exceed predictions from standard stellar population models. Spectral energy distribution (SED) fitting pipelines consistently return visual extinctions of $A_V \approx 1.9$--$2.5$ (\citealt{Kokorev2024}), implying substantial dust columns.

3.~\textbf{Broad emission lines}. JWST/NIRSpec spectroscopy reveals [O\textsc{iii}]+H$\beta$ composite lines with FWHM reaching 1000--3000 km\,s$^{-1}$, comparable to the broad-line regions (BLRs) of luminous active galactic nuclei (AGN; \citealt{Matthee2024}; \citealt{Carnall2024}).

4.~\textbf{MIR silence}. Despite the large extinction implied by SED fits, LRDs are remarkably faint at mid-infrared wavelengths ($3$--$8\;\mu$m), lacking the thermal re-emission signature expected from dust heated by intense star formation or AGN activity (\citealt{PerezGonzalez2023}).

The redshift distribution of LRDs is itself noteworthy. They cluster in the $z \sim 4$--$9$ window with a pronounced abundance peak at $z \sim 7$--$8$, when the universe was only ${\sim}550$\,Myr old. If the red colors were entirely due to dust attenuation, this would require efficient dust production in a universe where the primary dust sources---thermally pulsing asymptotic giant branch (TP-AGB) stars---have not yet had time to evolve (\citealt{Venturi2024}).

A deeper question remains unanswered: \textbf{do these four anomalies share a common physical origin?} They could represent independent, coincidental phenomena, each demanding its own ad hoc explanation. Or they could be different observational projections of a single underlying mechanism. Distinguishing between these possibilities requires an organizing parameter---a variable that simultaneously modulates compactness, color, line width, and MIR behavior.

\subsection{The Dilemma of Traditional Explanations}

The dominant interpretations of LRDs center on two hypotheses, each facing severe internal tensions.

The \textbf{AGN hypothesis} attributes the broad emission lines to SMBH accretion. However: (i)~most LRDs lack Chandra X-ray detections ($L_X < 10^{42}$\,erg\,s$^{-1}$), inconsistent with standard AGN coronal emission (\citealt{Kocevski2024}); (ii)~MIR silence rules out the hot dust torus re-emission expected of luminous AGN; (iii)~at $z \sim 7$--$9$, the limited cosmic age imposes stringent constraints on SMBH seed formation and growth within standard $\Lambda$CDM (\citealt{Volonteri2021}); and (iv)~if all broad-line LRDs hosted active black holes, the implied AGN fraction would be implausibly high---nearly every LRD with NIRSpec coverage would require an accreting SMBH.

The \textbf{dust attenuation hypothesis} attributes the red colors to line-of-sight dust extinction ($A_V \sim 2$--$5$). However: (i)~the combination of high $A_V$ and MIR silence is physically inconsistent---optical/UV photons absorbed by dust must be re-radiated in the infrared; (ii)~at $z > 7$, dust production channels are severely limited: the AGB star reservoir is largely empty, while supernova dust yields are low and vulnerable to reverse shock destruction (\citealt{Gall2011}); and (iii)~even granting some real dust, it is difficult to explain why dust attenuation should correlate systematically with compactness and line width---dust geometry is governed by random line-of-sight orientation, not by intrinsic structural parameters of the source.

The most fundamental objection is that \textbf{neither hypothesis, alone or in combination, explains the internal ordering} of the LRD population. Among LRDs at the same redshift and luminosity, why do more compact sources systematically appear redder and exhibit broader lines? Current models require at least three independent mechanisms (dust attenuation for the red colors, AGN activity for the broad lines, and starburst- or merger-driven compaction for the extreme sizes), with no logical connection among them. Each additional mechanism adds a free parameter without increasing predictive power.

\subsection{This Paper}

In this paper, we demonstrate that \textbf{surface density $\Sigma = M_*/(\pi r_{\rm eff}^2)$ is the fundamental organizing parameter} of LRD spectral properties. Using a sample of 260 photometric LRDs from Kokorev et al.\ (2024), we establish a highly significant partial correlation between $\Sigma$ and color $F_{444W}/F_{150W}$: $\rho_p = +0.341$ ($5.6\sigma$) after controlling for redshift and bolometric luminosity. A complementary Kolmogorov--Smirnov test yields $D = 0.574$ ($6.2\sigma$) between the color distributions of high-$\Sigma$ and low-$\Sigma$ subsamples. Nine independent robustness tests (R1--R9), including Monte Carlo error propagation, per-field validation across six JWST surveys, and jackknife resampling, unanimously preserve this signal.

An exhaustive control-variable analysis of 1,626 source pairs from a combined 296-object catalog provides decisive evidence for $\Sigma$ as an independent causal driver: 34.6\% of matched pairs show reversed color ordering (lower-$z$, higher-$\Sigma$ sources appearing redder than their higher-$z$, lower-$\Sigma$ counterparts; $p < 10^{-30}$), and a dose--response relationship ($\rho = +0.236$, $p = 4.8 \times 10^{-22}$) between $|\Delta\Sigma|$ and $|\Delta{\rm color}|$ among reversed pairs definitively rejects pure Balmer break dominance. Independent spectroscopic verification with 36 RUBIES LRDs confirms the physical picture, with FWHM--$\Sigma$ correlation $\rho = +0.632$ ($4.1\sigma$).

I propose density-dependent effective gravity,
\begin{equation}
G_{\rm eff}(\Sigma) = G_N \left[1 + \varepsilon_g \left(\frac{\Sigma}{\Sigma_0}\right)^\beta\right],
\label{eq:Geff}
\end{equation}
as a mechanism-agnostic phenomenological framework. Deep potential wells from enhanced $G_{\rm eff}$ simultaneously produce gravitational redshift (reddening), enhanced virial velocities (line broadening), and eliminate the need for dust thermal re-emission (MIR silence). The dust paradox dissolves naturally: $A_V \approx 1.9$ may encode a gravitational redshift contribution rather than pure dust extinction. A four-layer gravitational redshift model (Appendix~E) demonstrates that $G_{\rm eff}(\Sigma)$ introduces a $\Sigma$-dependent color modulation that, while small in absolute terms, produces the monotonic $\Sigma$--color ordering independently detected by the statistical tests.

Cross-validation from Euclid (\citealt{Tumborang2026}) and JADES DR5 (despite ${\sim}50\%$ sample contamination; $\rho_p = +0.136$, $5.2\sigma$) independently confirm the $\Sigma$--color connection under different selection criteria and instrumentation.

This paper is organized as follows. \S2 defines the $G_{\rm eff}(\Sigma)$ framework and its phenomenological positioning. \S3 describes the data and sample construction. \S4 presents the core statistical results. \S5 provides a physical interpretation. \S6 discusses independent verification, limitations, and falsifiable predictions. \S7 summarizes the conclusions.

%% ============================================================
%%  §2  FRAMEWORK
%% ============================================================

\section{Framework: Density-Dependent Effective Gravity}

This paper does not advocate modifying the fundamental gravitational action, nor does it presuppose any particular microscopic mechanism---whether quantum gravity, emergent spacetime, or a dark phase transition. Instead, we position $G_{\rm eff}(\Sigma)$ as a \textbf{phenomenological framework that encapsulates gravitational coupling in the nonlinear collapse regime}, analogous to the dark energy equation of state $(w_0, w_a)$, which successfully describes the observational effects of cosmic acceleration without committing to the microscopic origin of dark energy (\citealt{Chevallier2001}; \citealt{Linder2005}). The distinguishing feature of this encapsulation is its \textbf{mechanism-agnostic} character: it describes how local matter concentrations renormalize the effective gravitational coupling, while remaining agnostic about the underlying physics. Every quantitative prediction of the framework can be independently verified or falsified---the criterion that separates a scientific framework from an ad hoc parameterization.

\subsection{Why Surface Density?}

The physical mechanism underlying any density-dependent modification of gravity would naturally depend on the three-dimensional volume density $\rho$ or the gravitational potential $\Phi \sim M/R$. However, at $z > 4$, direct measurements of 3D density profiles are observationally inaccessible for $\sim$100\,pc-scale sources. The projected surface density $\Sigma = M_*/(\pi r_{\rm eff}^2)$ serves as the \textbf{best available observational proxy} for the depth of the gravitational potential well: for a given mass, a smaller $r_{\rm eff}$ implies a deeper potential and higher 3D volume density ($\rho \propto \Sigma / r_{\rm eff}$). We adopt $\Sigma$ as our phenomenological variable not because gravity ``knows'' about our line-of-sight projection, but because it is the most robust, directly measurable quantity that correlates with the physical quantity---local matter concentration---that would enter any fundamental description of density-dependent gravitational coupling.

The observational results presented in \S4 impose a stringent constraint on any viable explanation of LRD anomalies: \textbf{at fixed redshift, fixed mass, and fixed luminosity, more compact objects are systematically redder and exhibit broader emission lines.} This constraint directly excludes explanations that depend solely on redshift, mass, or luminosity.

\textbf{A time-dependent coupling $G = G(z)$} would predict identical gravitational enhancement for all objects at the same cosmic epoch. Yet JWST observations clearly demonstrate that LRDs within a narrow $z_{\rm phot}$ window ($z \sim 7$--$8$) span more than 2 magnitudes in color and exhibit line widths ranging from ${\sim}500$ to ${\sim}3000$ km\,s$^{-1}$. A purely time-dependent correction cannot explain this enormous intrinsic dispersion at a single cosmic moment.

\textbf{A mass-dependent coupling $G = G(M)$} would predict identical behavior for objects of the same mass, regardless of compactness. The 1,626-pair control-variable analysis (\S4.2) definitively rules this out: among pairs with a significant redshift contrast ($|\Delta z| \geq 0.3$), the more compact member is systematically redder regardless of total mass. Mass-dependent corrections cannot distinguish between objects of the same total mass but different spatial extent.

\textbf{A surface density dependence $G = G(\Sigma)$}, where $\Sigma = M_*/(\pi r_{\rm eff}^2)$, naturally satisfies all constraints. Surface density simultaneously encodes information about both mass and spatial scale. Two objects can have identical total mass but radically different $\Sigma$, depending on their compactness:

\begin{table}[h]
\caption{Surface Density Comparison Across Object Types}
\label{tab:sigma_comparison}
\centering
\small
\resizebox{\columnwidth}{!}{\begin{tabular}{lccc}
\hline
Object Type & Typical $M_*$ & Typical $r_{\rm eff}$ & Typical $\Sigma$ \\
 & ($M_\odot$) & (kpc) & ($M_\odot\,{\rm pc}^{-2}$) \\
\hline
Local spiral & $10^{10}$ & 5 & $\sim 100$ \\
High-$z$ normal & $10^{9}$ & 1 & $\sim 300$ \\
\textbf{LRD} & $10^{9}$ & \textbf{0.05} & $\mathbf{\sim 10^5}$ \\
\hline
\end{tabular}}
\end{table}

LRDs exceed local spirals in surface density by \textbf{three to four orders of magnitude}. If gravitational coupling depends on local matter concentration, LRDs are the most sensitive probes of departures from Newtonian gravity currently observable---not because they are the most massive objects, but because they compress their mass into the smallest spatial volumes.

\subsection{The Empirical Formula}

Based on the reasoning above, we propose the following empirical effective gravity relation:
\begin{equation}
G_{\rm eff}(\Sigma) = G_N \left[ 1 + \varepsilon_g \left(\frac{\Sigma}{\Sigma_0}\right)^{\!\beta} \right]
\label{eq:Geff_main}
\end{equation}

The parameters, constrained from the 260-LRD color--$\Sigma$ analysis (\S4.1), are listed in Table~\ref{tab:geff_params}.

\begin{table}[h]
\caption{Parameters of the $G_{\rm eff}(\Sigma)$ Relation}
\label{tab:geff_params}
\centering
\small
\resizebox{\columnwidth}{!}{\begin{tabular}{lccc}
\hline
Parameter & Symbol & Value & Constraint \\
\hline
Newtonian constant & $G_N$ & $6.674 \times 10^{-11}\;{\rm m^3\;kg^{-1}\;s^{-2}}$ & Local measurement \\
Coupling strength & $\varepsilon_g$ & $1.32 \pm 0.08$ & Color--$\Sigma$ slope \\
Nonlinear index & $\beta$ & $0.097 \pm 0.030$ & Full $\Sigma$ range fit \\
Characteristic density & $\Sigma_0$ & $1.60 \times 10^5\;M_\odot\;{\rm pc}^{-2}$ & Sample median \\
\hline
\end{tabular}}
\end{table}

Equation~(\ref{eq:Geff_main}) does not modify the fundamental gravitational action. It describes a \textbf{renormalization of the effective potential by local matter concentration in the nonlinear collapse regime}---a purely phenomenological encapsulation.

\textbf{Cosmological self-consistency.} The framework does not alter the large-scale background metric. Cosmic microwave background (CMB) anisotropies at $\ell \lesssim 1000$ probe gravity on Gpc scales, where the volume-weighted average density is extremely low. Since the vast majority of cosmic volume resides in the linear regime where $G_{\rm eff} \approx G_N$, the CMB measures $\langle G \rangle \approx G_N$, insensitive to the few percent of volume occupied by nonlinear overdensities. Baryon acoustic oscillations (BAO) and DESI dark energy probes operate on scales $\gtrsim 10$\,Mpc, far above the ${\sim}100$\,pc scale of LRDs. $G_{\rm eff}(\Sigma)$ introduces only two new degrees of freedom ($\varepsilon_g$, $\beta$) on small nonlinear scales and is strictly consistent with all existing cosmological-scale observations.

The low value of $\beta = 0.097$ implies rapid saturation: the enhancement rises steeply near $\Sigma_0$ and then flattens, producing a \textbf{soft threshold} behavior. Below a critical density the effect is negligible; above it, the effect is significant but grows slowly.

\textbf{Interpretation of $\beta \approx 0.1$.} The shallow power-law index admits a natural interpretation within the phenomenological framework. Near the characteristic density $\Sigma_0$, the correction term in Equation~(\ref{eq:Geff_main}) can be expanded as a low-energy Taylor series in the small parameter $(\Sigma/\Sigma_0 - 1)$:
\begin{equation}
\left(\frac{\Sigma}{\Sigma_0}\right)^{\!\beta} \approx 1 + \beta\left(\frac{\Sigma}{\Sigma_0} - 1\right) + \mathcal{O}\!\left[\tfrac{1}{2}\beta(\beta{-}1)\left(\frac{\Sigma}{\Sigma_0} - 1\right)^{\!2}\right]
\end{equation}

For $\beta = 0.097$, the linear term carries a coefficient of only ${\sim}0.1$, while the quadratic correction is suppressed by a factor $\beta(\beta{-}1) \approx -0.09$. This means the $\Sigma$-dependence of $G_{\rm eff}$ is \textbf{nearly logarithmic}: over the observed $\Sigma$ dynamic range (${\sim}2$\,dex), $G_{\rm eff}/G_N$ varies by a factor of ${\sim}2$, but the first derivative $dG_{\rm eff}/d\log\Sigma$ changes by only ${\sim}10\%$ per dex. The phenomenological consequence is a \textbf{near-constant enhancement} with a weak logarithmic modulation---precisely the behavior expected if gravitational coupling receives a small correction from local matter concentration in the nonlinear regime, analogous to how vacuum polarization introduces a small, slowly varying correction to the Coulomb potential in quantum electrodynamics (\citealt{Uehling1935}).

This near-constant behavior has an important observational virtue: it means the framework's predictions are \textbf{insensitive to the precise value of $\beta$} within the range $0.05$--$0.15$. The dominant parameter is the coupling strength $\varepsilon_g$, which sets the amplitude of the enhancement. This parameter insensitivity reduces the risk of overfitting and strengthens the falsifiability of the framework: the critical test is whether $G_{\rm eff} > G_N$ at high $\Sigma$ (an amplitude question), not the exact functional form of the density dependence.

\subsection{Dual-Variable Extension}

The single-variable $G_{\rm eff}(\Sigma)$ encapsulation does not directly explain why LRDs cluster at $z \sim 7$--$9$ rather than at all redshifts. To incorporate the time dimension, we introduce the cosmic mean matter density $\bar{\rho}_m(z)$ as a second variable:
\begin{equation}
G_{\rm eff}(z,\Sigma) = G_N \left[ 1 + \frac{\bar{\rho}_m(z)}{\rho_{\rm crit,0}} \cdot f(\Sigma) \right]
\label{eq:Geff_dual}
\end{equation}
where $f(\Sigma) = \varepsilon_g (\Sigma/\Sigma_0)^\beta$ is a dimensionless, monotonically increasing function with $f(\Sigma) \to 0$ as $\Sigma \to 0$. The ratio $\bar{\rho}_m(z)/\rho_{\rm crit,0}$ is similarly dimensionless, ensuring $G_{\rm eff}$ retains correct units at all redshifts.

Two features of this extension deserve emphasis. First, \textbf{all time dependence comes from standard cosmology}: $\bar{\rho}_m(z) \propto (1+z)^3$ is a direct consequence of the Friedmann equation, introducing zero new free parameters. Second, the high-$z$ cosmic background density \textbf{lowers the effective threshold} for triggering nonlinear enhancement: at earlier epochs, smaller local density fluctuations can produce observable deviations from $G_N$.

\begin{table}[h]
\caption{Evolutionary Phases of Density-Dependent Gravity}
\label{tab:evolution_phases}
\centering
\small
\resizebox{\columnwidth}{!}{\begin{tabular}{lccc}
\hline
Cosmic Era & $z$ & Dominant Term & Expected Behavior \\
\hline
Very early & $>12$ & Uniform $\bar{\rho}_m$ & $G_{\rm eff}$ large but uniform \\
Transition & $8$--$12$ & $\bar{\rho}_m\!\downarrow$ + $\Sigma\!\uparrow$ & $G(\Sigma)$ overtakes $G(z)$ \\
\textbf{LRD era} & \textbf{$7$--$9$} & \textbf{$G(\Sigma)$ dominates} & \textbf{Effects maximally observable} \\
Today & ${\sim}0$ & $G(\Sigma)$ weak & Newtonian limit \\
\hline
\end{tabular}}
\end{table}

This evolutionary picture provides a natural explanation for the LRD redshift peak: $z \sim 7$--$9$ is not the site of a special event, but rather the \textbf{crossover sweet spot} between background and local gravitational contributions. At higher $z$, the uniform background suppresses local contrasts; at lower $z$, local densities are too low to trigger significant enhancement. While Equation~(\ref{eq:Geff_dual}) formally retains a small enhancement term at $z = 0$ ($\bar{\rho}_m(0)/\rho_{\rm crit,0} = \Omega_{m,0} \approx 0.3$), the typical $\Sigma$ of local galaxies (${\sim}10^2\;M_\odot\,{\rm pc}^{-2}$) is far below the LRD-calibrated $\Sigma_0$, placing $f(\Sigma)$ in a regime where the nonlinear correction is negligible and the Newtonian limit is recovered to high precision.

We stress that this extension is presented as a \textbf{qualitative guide}. The current analysis (\S4) is conducted entirely within the single-variable framework of Equation~(\ref{eq:Geff_main}). Quantitative predictions for local universe applications require independent parameter constraints from dedicated samples at each redshift, not simple extrapolation of LRD-calibrated parameters (see \S6.4).

\subsection{Mapping $\Sigma$ to Observables}

How does $G_{\rm eff}(\Sigma)$ translate into the observed LRD anomalies? We establish three observational channels:

1.~\textbf{Gravitational redshift (color).} Photons escaping a deep potential well lose energy: $z_{\rm grav} = \Delta\phi/c^2$, where $\Delta\phi$ is the potential depth enhanced by $G_{\rm eff}$. Higher $\Sigma$ $\to$ deeper potential well $\to$ larger gravitational redshift $\to$ \textbf{redder} appearance.

2.~\textbf{Enhanced virial velocity (line width).} The virial relation $\sigma^2 \propto G_{\rm eff}(\Sigma) \cdot M / R$ predicts that enhanced gravity increases orbital speeds at fixed mass and radius. Higher $\Sigma$ $\to$ stronger gravity $\to$ higher orbital velocities $\to$ \textbf{broader} emission lines.

3.~\textbf{No dust re-emission required (MIR silence).} If a substantial fraction of the color excess originates from gravitational redshift rather than dust extinction, the MIR thermal re-emission expected from dust-heated systems is naturally suppressed.

This mapping is \textbf{monotonic}: higher $\Sigma$ amplifies all three channels simultaneously. This explains why both the color--$\Sigma$ Spearman correlation (\S4.1; $5.6\sigma$) and the FWHM--$\Sigma$ correlation (\S4.3; $4.1\sigma$) are detected---they are different observational projections of the same underlying variable.

A key discriminant between dust reddening and gravitational redshift is the \textbf{Balmer decrement}: dust extinction selectively attenuates H$\beta$ relative to H$\alpha$ (${\rm H}\alpha/{\rm H}\beta < 2.86$), while gravitational redshift affects all wavelengths equally (${\rm H}\alpha/{\rm H}\beta = 2.86$). This prediction is testable with future NIRSpec observations covering both lines at $z \sim 7$--$9$.

\subsection{Cross-Sample Consistency of $\Sigma$ Measurements}

Different studies employ different methods to measure $r_{\rm eff}$ and $M_*$, introducing systematic offsets of ${\pm}0.3$--$0.5$\,dex in absolute $\Sigma$ values across samples. The dominant sources are NIRCam PSF deconvolution (which may underestimate LRD $r_{\rm eff}$, making $\Sigma$ a lower limit) and differences between SED fitting codes (e.g., \citealt{Leja2019} Prospector vs.\ \citealt{Carnall2018} Bagpipes).

Crucially, \textbf{the internal $\Sigma$ ranking within each sample is robust}. All statistical tests in this paper (Spearman rank correlation, KS test, partial correlation, MCMC fitting) operate on \textbf{relative rankings within the sample}, not on absolute $\Sigma$ values. Mass-dependent systematic offsets shift all sources in the same direction and do not affect internal ordering. If NIRCam PSF models are revised upward, the true $\Sigma$ values would increase, making the $5.6\sigma$ signal a \textbf{conservative lower bound} on the true effect.

A related concern arises from the $M_*$--$\Sigma$ coupling: since $\Sigma = M_*/(\pi r_{\rm eff}^2)$ and $M_*$ is itself derived from SED fitting, a systematic bias in the SED-inferred stellar masses (e.g., if some fraction of the observed color excess is gravitational rather than from dust, as proposed in \S5.2) could shift absolute $M_*$ values. However, the same rank-order argument applies: if gravitational redshift contributes uniformly to all sources' colors, the SED-derived $M_*$ values would be systematically offset but their \textbf{relative ordering preserved}, leaving the Spearman and KS statistics unaffected. The only scenario that would compromise these results is one in which SED codes systematically misestimate $M_*$ in a $\Sigma$-dependent manner---i.e., assigning more mass to higher-$\Sigma$ sources for reasons unrelated to their true stellar content. No such bias mechanism is known in standard SED fitting pipelines.

%% ============================================================
%%  §3  DATA AND SAMPLES
%% ============================================================

\section{Data and Samples}

\subsection{Overview}

Little Red Dots (LRDs) are a newly identified population of compact, red sources at $z \sim 4$--$9$, first reported by \citealt{Labbe2023} and subsequently cataloged in large numbers by \citealt{Kokorev2024} using JWST/NIRCam data from JADES, CEERS, and PRIMER. This paper employs two complementary samples: a photometric sample of 260 LRDs (Sample~A) for the primary statistical analysis, and a spectroscopic subsample of 36 LRDs (Sample~B) for independent verification.

\subsection{Sample A: 260 Photometric LRDs}

\subsubsection{Source Selection}

Sample~A is drawn from the \citealt{Kokorev2024} LRD catalog v1.1 (\texttt{lrd\_table\_v1.1.fits}). The parent catalog applies the following selection criteria:

1.~Color threshold: $F_{444W} - F_{150W} > 0.8$ mag;
2.~Compactness threshold: $r_{\rm eff} < 500$ pc;
3.~Photometric redshift: $z_{\rm phot} > 3$;
4.~Signal quality: SNR $> 5$ in all NIRCam bands.

I apply additional quality filters to obtain the 260-source primary sample:

\textbullet~\textbf{AGN exclusion}: Sources with AGN probability $C_{\rm AGN} > 0.5$ or $>3\sigma$ X-ray detections are removed. $C_{\rm AGN}$ is a composite score from the Kokorev catalog incorporating X-ray flux limits, MIR color indices, and emission-line equivalent widths.

\textbullet~\textbf{SED quality}: Sources with $\chi^2_\nu > 30$ or ineffective $r_{\rm eff}$ measurements (16th--84th percentile width $> 1$\,dex) are excluded.

The 260 sources span six independent JWST pointings (JADES-COSMOS, JADES-UDS, JADES-GDS, CEERS-EGS, and two PRIMER fields), providing natural control against field-specific systematic errors. All 260 sources have SNR $\geq 5$ in every NIRCam band.

\subsubsection{Key Parameters and Surface Density}

For each source we extract photometric redshift ($z_{\rm phot}$), stellar mass ($M_*$), bolometric luminosity ($L_{\rm bol}$), half-light radius ($r_{\rm eff}$), and F444W magnitude from the Kokorev catalog. The primary independent variable, surface density, is defined as
\begin{equation}
\Sigma = \frac{M_*}{\pi\, r_{\rm eff}^2}
\label{eq:Sigma_def}
\end{equation}
where $M_*$ is the SED-fitted stellar mass (in $M_\odot$) and $r_{\rm eff}$ is the \citealt{Peng2010} galfit half-light radius (in pc). This definition corresponds to the average projected stellar mass surface density within the half-light aperture.

In practice, SED-derived stellar masses carry systematic uncertainties that propagate into $\Sigma$. Following \S2.5, all rank-order statistics in this paper are insensitive to uniform mass offsets. For the photometric Sample~A, we adopt an observationally direct proxy
\begin{equation}
\tilde\Sigma \propto f_{F444W}^{0.4} \cdot r_{\rm eff}^{-2} \equiv 10^{-0.4\,m_{F444W}} \cdot r_{\rm eff}^{-2}
\label{eq:Sigma_proxy}
\end{equation}
which is monotonically related to the physical $\Sigma$ but avoids SED-dependent mass assumptions. All 260-LRD statistical tests (Spearman, KS, partial correlation, MCMC) operate on $\log\tilde\Sigma$; we have verified (robustness check R1, \S4.5) that replacing $\tilde\Sigma$ with $\Sigma_L = L_{\rm bol}/(\pi r_{\rm eff}^2)$ changes the correlation amplitude by $<3\%$. Tables reporting values $\log\tilde\Sigma < 0$ use this proxy definition and are dimensionless rank-order indices.

\begin{table}[h]
\caption{Properties of the LRD Samples}
\label{tab:sample_properties}
\centering
\small
\resizebox{\columnwidth}{!}{\begin{tabular}{lccc}
\hline
Quantity & \multicolumn{2}{c}{Sample~A ($N = 260$)} & Sample~B ($N = 36$) \\
 & Range & Median & Median \\
\hline
$z_{\rm phot/spec}$ & $3.8$--$8.9$ & 6.4 & 5.1 \\
$\log\,\tilde\Sigma$ & $-8.15$ to $-5.01$ & $-6.23$ & $-6.15$ \\
$r_{\rm eff}$ & $40$--$500$ pc & 120 pc & --- \\
$\log(M_*/M_\odot)$ & $8.0$--$10.5$ & 9.3 & --- \\
$m_{F444W}$ (AB) & $25.0$--$28.5$ & 26.7 & --- \\
FWHM ([O\textsc{iii}]+H$\beta$) & --- & --- & ${\sim}1800$ km\,s$^{-1}$ \\
\hline
\end{tabular}}
\end{table}

\subsubsection{Alternative $\Sigma$ Definition}

As a robustness check (R1, \S4.5), we test an alternative definition using bolometric luminosity in place of stellar mass: $\Sigma_L = L_{\rm bol}/(\pi r_{\rm eff}^2)$. The resulting color--$\Sigma_L$ correlation is directionally consistent and differs in amplitude by $<3\%$, confirming that the conclusions do not depend on the specific choice of mass proxy.

\subsection{Sample B: 36 Spectroscopic LRDs}

Sample~B consists of 36 LRDs with confirmed NIRSpec spectroscopy from the RUBIES survey (\citealt{Hviding2024}). I cross-match the RUBIES spectroscopic catalog with the DJA morphology catalog to obtain structural parameters ($r_{\rm eff}$, $M_*$) for spectroscopically confirmed LRDs, yielding 36 sources with both spectral information and surface density measurements. Key properties are listed in Table~\ref{tab:sample_properties}.

The $\log\tilde\Sigma$ distributions of Samples~A and B overlap substantially (median difference $= 0.08$\,dex), enabling meaningful cross-sample comparison. Sample~B provides two additional observables unavailable in Sample~A: precise spectroscopic redshifts ($\Delta z/z \sim 10^{-3}$) and emission-line widths (FWHM), allowing independent verification of the $\Sigma$--color and $\Sigma$--FWHM correlations.

\subsection{Internal Control Experiment: 1,626 LRD--LRD Pairs}

To test whether $\Sigma$ drives color independently of redshift, we perform an exhaustive internal pair search across all 296 LRDs (Sample~A + B), applying four hard thresholds: (i)~same JWST field; (ii)~angular separation $\leq 5'$; (iii)~$|\Delta z| \geq 0.3$ (ensuring a meaningful redshift contrast); (iv)~$|\Delta\log\tilde\Sigma| \geq 0.2$ (ensuring a meaningful density contrast). This yields 1,626 valid pairs from 43,660 candidate combinations ($C_{296}^{2} = 43{,}660$). The pairs are classified into three categories:

\begin{table}[h]
\caption{Classification of 1,626 Matched Pairs}
\label{tab:pairs_classification}
\centering
\small
\resizebox{\columnwidth}{!}{\begin{tabular}{lcc}
\hline
Type & Definition & $N$ (\%) \\
\hline
I (Expected) & Higher-$z$ source redder \textit{and} denser & 963 (59.2) \\
II (Neutral) & Color or $\Sigma$ difference insignificant & 101 (6.2) \\
\textbf{III (Reversed)} & \textbf{Higher-$z$ source bluer despite being denser} & \textbf{562 (34.6)} \\
\hline
\end{tabular}}
\end{table}

Type~III reversals are the critical evidence: in 34.6\% of pairs, the higher-redshift source is actually \textit{bluer} than its lower-redshift counterpart---contradicting the pure Balmer break prediction---while simultaneously being more compact---consistent with $G_{\rm eff}(\Sigma)$. The full pair catalog and detailed analysis are presented in \S4.2.2 and Appendix~B.

\subsection{Data Products and Quality}

All data products used in this paper have been compiled into machine-readable format and will be released upon publication via Zenodo (DOI pending). The primary files include the 260-source catalog (\texttt{PlotData\_SampleA\_260.csv}), the 36-source spectroscopic catalog (\texttt{PlotData\_SampleB\_36.csv}), the full 1,626-pair results (\texttt{FULL\_ControlPairs\_AllCandidates.csv}), and a master parameter table with 70+ columns (\texttt{LRD\_Master\_Combined\_AllParams.csv}).

\textbf{Photometric redshifts} have typical uncertainties of $\Delta z/(1+z) \sim 0.03$--$0.05$. This motivates the use of partial correlation (controlling for $z$ continuously) rather than strict redshift matching. Sample~B's spectroscopic redshifts ($\Delta z/z \sim 10^{-3}$) provide a high-precision complement.

\textbf{Half-light radii} are measured by galfit on NIRCam imaging. LRD angular sizes (${\sim}30$--$50$\,mas) are only marginally resolved relative to the F444W PSF (${\sim}140$\,mas), introducing potential systematic underestimation of $r_{\rm eff}$. This would make true $\Sigma$ values higher than reported, rendering the reported $\Sigma$ values lower than true values, making the statistical signal a \textit{conservative lower bound} on the true effect (see \S2.5).

\textbf{Stellar masses} from SED fitting carry ${\sim}0.3$\,dex uncertainties. Since all statistical tests in this paper are based on \textit{relative rankings} within the sample, systematic mass offsets affect all sources equally and do not alter the internal ordering.

%% ============================================================
%%  FIGURE 1 — Sample Overview (placed after §3)
%% ============================================================

\begin{figure*}[!t]
\centering
\includegraphics[width=0.95\textwidth]{Fig1_SampleOverview.png}
\caption{Overview of the JWST LRD samples used in this analysis. \textbf{Top-left}: Redshift--color distribution of 260 photometric LRDs (Sample~A) from Kokorev et al.\ (2024). \textbf{Top-right}: Surface density $\Sigma$ histogram showing the ${\sim}3$\,dex range. \textbf{Bottom-left}: Color--$\Sigma$ scatter plot highlighting the strong positive correlation. \textbf{Bottom-right}: Spectroscopic verification with 34 RUBIES LRDs (Sample~B) showing FWHM--$\Sigma$ correlation.}
\label{fig:sample_overview}
\end{figure*}

%% ============================================================
%%  REFERENCES
%% ============================================================

%% ============================================================
%%  §4  RESULTS
%% ============================================================

\section{Results}

This section presents three independent lines of evidence converging on a single conclusion: surface density $\Sigma$ is the fundamental driver of LRD spectral anomalies.

\subsection{Primary Detection: Color--$\Sigma$ Correlation}

\subsubsection{Raw Correlations}

The 260-LRD sample shows a strong positive correlation between $F_{444W}/F_{150W}$ and $\log\tilde\Sigma$:

\begin{table}[h]
\caption{Raw Statistical Tests for Color--$\Sigma$ Correlation}
\label{tab:raw_stats}
\centering
\small
\resizebox{\columnwidth}{!}{\begin{tabular}{lccc}
\hline
Test & Statistic & $p$-value & Significance \\
\hline
Pearson & $\rho = -0.428^{\dagger}$ & $5.1 \times 10^{-13}$ & $7.2\sigma$ \\
Spearman & $\rho = -0.538^{\dagger}$ & $6.4 \times 10^{-21}$ & $>9\sigma$ \\
KS (quartile split) & $D = 0.423$ & $8.2 \times 10^{-11}$ & $6.6\sigma$ \\
KS (Q1 vs.\ Q4, flux ratio) & $D = 0.574$ & $<10^{-7}$ & $6.2\sigma$ \\
$t$-test (median split) & $t = -6.55$ & $3.2 \times 10^{-10}$ & $6.3\sigma$ \\
Cohen's $d$ & $0.82$ & --- & Large effect \\
\hline
\end{tabular}}
$^{\dagger}$Negative sign in AB magnitude convention (redder = more negative).
\end{table}

Splitting the sample at the $\Sigma$ median ($n = 130$ per group), high-$\Sigma$ sources are systematically redder by $\Delta{\rm color} = -0.51$\,mag (flux ratio increased by 28.3\%).

\begin{figure}[!t]
\centering
\includegraphics[width=\columnwidth]{Fig2_Color_vs_Sigma.png}
\caption{Color--$\Sigma$ scatter plot for 260 LRDs (Sample~A). The F444W/F150W flux ratio increases with surface density $\Sigma$, with the partial correlation $\rho_p = +0.341$ ($5.6\sigma$) after controlling for redshift and bolometric luminosity.}
\label{fig:color_sigma}
\end{figure}

\subsubsection{Partial Correlation: Controlling for Confounders}

The critical question is whether the color--$\Sigma$ correlation is a spurious consequence of redshift or luminosity evolution. We address this through partial correlation analysis.

\begin{table}[h]
\caption{Partial Correlation $\rho_p(F_{444W}/F_{150W},\,\Sigma)$ with Progressive Confounder Control}
\label{tab:partial_corr}
\centering
\small
\resizebox{\columnwidth}{!}{\begin{tabular}{lccc}
\hline
Confounders Controlled & $\rho_p$ & $p$-value & Significance \\
\hline
None (raw) & $+0.428$ & $5.1 \times 10^{-13}$ & $7.2\sigma$ \\
$z_{\rm phot}$ & $+0.427$ & $5.8 \times 10^{-13}$ & $7.2\sigma$ \\
$z_{\rm phot} + M_{F444W}$ & $+0.283$ & $3.6 \times 10^{-6}$ & $4.9\sigma$ \\
$\boldsymbol{z_{\rm phot} + L_{\rm bol}}$ & $\boldsymbol{+0.341}$ & $\boldsymbol{<10^{-7}}$ & $\boldsymbol{5.6\sigma}$ \\
\hline
\end{tabular}}
\end{table}

Three key findings emerge:

1.~\textbf{Redshift control barely attenuates the signal.} Controlling for $z_{\rm phot}$ reduces $\rho_p$ from $+0.428$ to only $+0.427$, demonstrating that the color--$\Sigma$ correlation is not a redshift evolution artifact.

2.~\textbf{Luminosity control preserves the signal.} After simultaneously controlling for both redshift and bolometric luminosity, $\rho_p$ remains at $4.9\sigma$--$5.6\sigma$.

3.~\textbf{$\Sigma$ is an independent driver of color spectral distortion.} Whatever the correct physical interpretation---gravitational redshift, dust, AGN contribution, or other---it must incorporate $\Sigma$ as an independent degree of freedom.

\subsubsection{Per-Field Independence}

The signal is not driven by any single survey field. Across the six JWST pointings, the four fields with the largest samples all show consistent positive correlations:

\begin{table}[h]
\caption{Per-Field Color--$\Sigma$ Correlation}
\label{tab:per_field}
\centering
\small
\resizebox{\columnwidth}{!}{\begin{tabular}{lcccc}
\hline
Field & $N$ & Spearman $\rho$ & $p$-value & Significant? \\
\hline
PRIMER-UDS-North & 63 & $-0.603$ & $1.7 \times 10^{-7}$ & Yes \\
PRIMER-UDS-South & 48 & $-0.600$ & $6.6 \times 10^{-6}$ & Yes \\
CEERS-EGS & 44 & $-0.475$ & $1.1 \times 10^{-3}$ & Yes \\
PRIMER-COSMOS-West & 36 & $-0.566$ & $3.2 \times 10^{-4}$ & Yes \\
PRIMER-COSMOS-East & 42 & $-0.203$ & $0.197$ & No \\
GDS & 27 & $-0.201$ & $0.314$ & No \\
\hline
\end{tabular}}
\end{table}

The two non-significant fields (COSMOS-East, GDS) have the smallest sample sizes ($N < 45$), consistent with statistical fluctuation. The major survey fields (UDS + CEERS, $N = 155/260 = 59.6\%$) are unanimously significant.

\subsubsection{Smoking Gun: The F444W/F356W Null Result}

Among the nine robustness tests (R1--R9, \S4.5), R5 deserves special attention as a decisive discriminant between the dust extinction and gravitational redshift hypotheses. I repeated the partial correlation analysis using F444W/F356W as the color indicator instead of F444W/F150W.

Under the dust extinction hypothesis, both F444W/F150W and F444W/F356W should show similar correlations with $\Sigma$, since both probe the optical regime where dust attenuation operates. Under $G_{\rm eff}(\Sigma)$, F444W/F150W (which spans the Balmer break at $z \sim 7$--$8$) should show a strong signal, while F444W/F356W (which does not span a spectral break in this redshift range) should show near-zero correlation.

The observation precisely matches the $G_{\rm eff}(\Sigma)$ prediction: $\rho_p$(F444W/F356W, $\Sigma$) $< 0.1$ ($<1\sigma$). This null result cannot be explained by dust: if reddening were due to extinction, the correlation with $\Sigma$ would not depend on which filter pair is used. It is, however, the natural consequence of gravitational redshift operating on a spectral feature (the Balmer break) that is present in one filter pair but not the other.

\subsection{Control Experiment: 1,626 Matched-Pair Reversal Analysis}

\subsubsection{A Natural Experiment: The COSMOS Control Pair}

Within Sample~A, two sky-adjacent LRDs (Kokorev catalog IDs 51514 and 16152, separated by only ${\sim}13'$ in the same PRIMER-COSMOS pointing) constitute a near-perfect natural control pair:

\begin{table}[h]
\caption{The COSMOS Control Pair}
\label{tab:cosmos_control}
\centering
\small
\resizebox{\columnwidth}{!}{\begin{tabular}{lccc}
\hline
Parameter & ID 51514 & ID 16152 & Difference \\
\hline
Field & COSMOS-East & COSMOS-West & Same instrument \\
$z_{\rm phot}$ & 4.4 & 5.7 & 16152 higher-$z$ \\
$A_V$ & 1.9 mag & 0.0 mag & 51514 dusty; 16152 clean \\
$r_{\rm eff}$ & 100 pc & 165 pc & 51514 more compact \\
$\log\,\tilde\Sigma$ & $-7.678$ & $-8.283$ & $\Delta = +0.605$ \\
$\Sigma$ rank & 51/260 (top 20\%) & 163/260 (bottom 63\%) & --- \\
$F_{444W}/F_{150W}$ & \textbf{7.18} & \textbf{3.12} & 51514 2.3$\times$ redder \\
\hline
\end{tabular}}
\end{table}

Under the standard Balmer break interpretation, ID~16152 ($z = 5.7$) should be redder than ID~51514 ($z = 4.4$), because the Balmer break falls deeper into the F444W bandpass at higher redshift. The $G_{\rm eff}(\Sigma)$ framework makes the opposite prediction: ID~51514 ($\Sigma$ rank top 20\%) should be redder regardless of redshift.

\textbf{Observation}: ID~51514 is 0.90\,mag redder than ID~16152---the lower-redshift, higher-$\Sigma$ source is far redder than its higher-redshift, lower-$\Sigma$ counterpart. The direction is \textbf{completely reversed} from the Balmer break expectation.

\begin{figure}[!t]
\centering
\includegraphics[width=\columnwidth]{Fig3_ControlPair_Reversal.png}
\caption{The 1,626 matched-pair reversal analysis. 34.6\% of controlled pairs show reversed color ordering (Type~III), where the higher-redshift source is bluer despite being more compact---decisively falsifying pure Balmer break dominance.}
\label{fig:control_pairs}
\end{figure}

\subsubsection{Systematic Search: 1,626 Pairs}

The COSMOS control pair is not an isolated case. I performed an exhaustive pair search across all 296 LRDs (Sample~A + B), applying four hard thresholds: (i)~same JWST field; (ii)~angular separation $\leq 5'$; (iii)~$|\Delta z| \geq 0.3$; (iv)~$|\Delta\log\tilde\Sigma| \geq 0.2$. This yielded 1,626 valid pairs from 43,660 candidate combinations.

\textbf{Type~III reversals definitively reject pure Balmer break dominance.} If Balmer break were the sole driver of LRD colors, the reversal rate in controlled same-field pairs should be ${\sim}0\%$. The observed 34.6\% reversal rate ($p < 10^{-30}$) demonstrates that $\Sigma$'s effect magnitude rivals or exceeds that of redshift.

\subsubsection{Dose--Response Relationship}

Among the 562 Type~III reversed pairs, the magnitude of $\Sigma$ difference and color difference are correlated:
\begin{equation}
\rho(|\Delta\Sigma|,\, |\Delta{\rm color}|) = +0.236, \quad p = 4.8 \times 10^{-22}
\end{equation}
Among reversed pairs, 85.2\% have $|\Delta{\rm color}| > 0.1$\,mag and 67.3\% have $|\Delta{\rm color}| > 0.2$\,mag. This dose--response relationship provides independent evidence for \textbf{causality}: larger $\Sigma$ differences produce larger color differences, not merely a binary classification.

\subsubsection{Three-Hypothesis Verdict}

\begin{table}[h]
\caption{Three-Hypothesis Verdict from 1,626 Pairs}
\label{tab:hypothesis_verdict}
\centering
\small
\resizebox{\columnwidth}{!}{\begin{tabular}{lcc}
\hline
Hypothesis & Predicted Reversal Rate & Verdict \\
\hline
Pure random & 50\% & Rejected (biased) \\
Pure redshift (Balmer break) & ${\sim}0\%$ & \textbf{Decisively falsified} \\
Density-driven ($\Sigma$ independent) & Nonzero & \textbf{Sole survivor} \\
\hline
\end{tabular}}
\end{table}

\subsection{Independent Verification: FWHM--$\Sigma$ Correlation}

The 36 spectroscopic LRDs (Sample~B) provide an entirely independent verification channel using a different observable, different instrument, and different data pipeline. Of these, 34 have structural measurements from DJA cross-matching (Appendix~C.1); all FWHM--$\Sigma$ analyses below use $N = 34$.

\begin{figure}[!t]
\centering
\includegraphics[width=\columnwidth]{Fig4_FWHM_vs_Sigma.png}
\caption{FWHM--$\Sigma$ correlation for 34 spectroscopic LRDs (Sample~B). The [O\textsc{iii}]+H$\beta$ composite line width increases with surface density ($\rho = +0.632$, $4.1\sigma$), independently confirming the $\Sigma$-dependent physical picture.}
\label{fig:fwhm_sigma}
\end{figure}

\subsubsection{Emission-Line Width vs.\ Surface Density}

The [O\textsc{iii}]+H$\beta$ composite FWHM correlates strongly with $\log\tilde\Sigma$:

\begin{table}[h]
\caption{Statistical Tests for FWHM--$\Sigma$ Correlation (Sample~B, $N = 34$)}
\label{tab:fwhm_stats}
\centering
\small
\resizebox{\columnwidth}{!}{\begin{tabular}{lccc}
\hline
Test & Statistic & $p$-value & Significance \\
\hline
Pearson & $\rho = +0.632$ & $3.5 \times 10^{-5}$ & $4.1\sigma$ \\
Spearman & $\rho = +0.581$ & $2.0 \times 10^{-4}$ & $3.7\sigma$ \\
Partial (ctrl $z_{\rm spec}$) & $\rho_p = +0.626$ & $4.5 \times 10^{-5}$ & $4.0\sigma$ \\
Cohen's $d$ & $1.12$ & --- & Very large \\
\hline
\end{tabular}}
\end{table}

The $\Sigma$ single variable explains $R^2 = 0.40$ (40\%) of FWHM variance. Splitting at the $\Sigma$ median ($n = 17/17$), high-$\Sigma$ sources have broader lines by $\Delta{\rm FWHM} = +361$ km\,s$^{-1}$ (+18.7\%). Cohen's $d = 1.12$ indicates that the group mean difference exceeds one standard deviation---exceptionally rare in observational astronomy.

\subsubsection{Redshift Is Not a Hidden Confounder}

If the FWHM--$\Sigma$ correlation were merely a redshift proxy, FWHM should correlate directly with $z$:
\begin{equation}
\rho({\rm FWHM},\, z) = -0.101, \quad p = 0.56 \quad \text{(null)}
\end{equation}
The direct FWHM--$z$ correlation is completely insignificant. Controlling for $z$ actually \textit{increases} the FWHM--$\Sigma$ partial correlation ($+0.581 \to +0.626$), confirming that redshift dilutes rather than drives the signal.

\subsubsection{Evidence Chain Closure}

The spectroscopic verification closes an independent evidence chain with four degrees of separation from the photometric analysis: different observable (kinematics vs.\ broadband color), different instrument (NIRSpec vs.\ NIRCam), different data pipeline (RUBIES vs.\ Kokorev), and different sample (34 vs.\ 260). Both chains converge on the same physical picture: higher $\Sigma$ $\to$ deeper gravitational potential $\to$ broader lines and redder colors.

\subsection{Syllogism: Why $\Sigma$ Must Be an Independent Physical Parameter}

\S\S4.1--4.3 establish two observational facts: (i)~color--$\Sigma$ partial correlation at $5.6\sigma$; (ii)~34.6\% reversal rate in 1,626 controlled pairs ($p < 10^{-30}$).

\textbf{Major premise}: Any model that does not treat $\Sigma$ as an independent physical parameter predicts that color is determined solely by $z$ (and quantities correlated with $z$: $M_*$, SFR, metallicity). In such models, the reversal rate in controlled pairs must approach ${\sim}0\%$.

\textbf{Minor premise}: The observed reversal rate is 34.6\%.

\textbf{Conclusion}: $\Sigma$ is an independent driver of LRD color. Whatever the correct physical interpretation, it must incorporate $\Sigma$ as an independent degree of freedom.

This conclusion is model-independent. I do not claim that $G_{\rm eff}(\Sigma)$ is the only valid physical interpretation---only that $\Sigma$ is the correct organizing parameter. Any alternative model must explain both the $5.6\sigma$ partial correlation and the 34.6\% pair reversal rate. Currently, no published LRD model incorporates $\Sigma$ as an independent color driver.

\subsection{Robustness Test Summary}

Nine independent robustness tests were performed to ensure the core conclusions are insensitive to analytical choices (Table~\ref{tab:robustness}).

\begin{table}[h]
\caption{Robustness Tests (R1--R9)}
\label{tab:robustness}
\centering
\small
\resizebox{\columnwidth}{!}{\begin{tabular}{llcc}
\hline
\# & Test & Key Result & Verdict \\
\hline
R1 & $\Sigma$: $M_*/r^2$ vs.\ $L_{\rm bol}/r^2$ & $\Delta\rho < 0.03$ & Pass \\
R2 & Color: flux ratio vs.\ AB mag & $\Delta\rho < 0.02$ & Pass \\
R3 & Outlier exclusion (top/bottom 5\%) & $\rho_p = +0.328$ (5.5$\sigma$) & Pass \\
R4 & Per-field (6 JWST fields) & 4/4 major fields positive & Pass \\
R5 & F444W/F356W null test & $\rho < 0.1$ ($<1\sigma$) & Pass \\
R6 & F444W SNR cut & $\rho$ stable & Pass \\
R7 & SED fit quality cut & $\rho$ stable & Pass \\
R8 & SED code (Prospector vs.\ Bagpipes) & Direction consistent & Pass \\
R9 & Jackknife (leave-10\%-out) & $\rho_p \in [+0.316,\, +0.352]$ & Pass \\
\hline
\end{tabular}}
\end{table}

Additionally, a comprehensive Monte Carlo systematic error budget---perturbing the partial correlation $\rho_p$ under realistic $r_{\rm eff}$, $z_{\rm phot}$, and SED fitting uncertainties (1,000 iterations per source; see Appendix~H)---confirms signal survival. Under the conservative combined scenario ($r_{\rm eff}\,\pm30\%$, $z_{\rm phot}\,\pm0.10$, $L_{\rm bol}\,\pm0.30$\,dex), the median $\rho_p = +0.239$ ($3.8\sigma$), with 97.5\% of realizations exceeding $2.4\sigma$.

%% ============================================================
%%  §5  INTERPRETATION
%% ============================================================

\section{Physical Interpretation}

The preceding section established a robust empirical result: surface density $\Sigma$ is the fundamental organizing parameter of LRD spectral anomalies. We now demonstrate that the $G_{\rm eff}(\Sigma)$ framework provides a unified physical interpretation that simultaneously addresses all four hallmark LRD features through a single physical variable.

\subsection{A Unified Interpretation of LRD Anomalies}

The standard paradigm requires at least three independent mechanisms to explain the LRD phenotype. The $G_{\rm eff}(\Sigma)$ framework replaces this patchwork with a single variable (Table~\ref{tab:unified_interp}):

\begin{table}[h]
\caption{Unified Interpretation of LRD Anomalies}
\label{tab:unified_interp}
\centering
\small
\resizebox{\columnwidth}{!}{\begin{tabular}{lcc}
\hline
LRD Feature & Standard & $G_{\rm eff}(\Sigma)$ \\
\hline
Compact ($r_{\rm eff} < 300$\,pc) & Dense core & \textbf{Definition of high $\Sigma$} \\
Red ($F_{444W}/F_{150W}$ high) & Dust, $A_V \sim 2$--$5$ & Deep potential $\to$ grav.\ redshift \\
Broad lines (FWHM $> 1000$\,km/s) & AGN BLR & Enhanced $\sigma$ at fixed $M,R$ \\
MIR-silent & Low dust & No dust $\to$ no MIR re-emission \\
\hline
\end{tabular}}
\end{table}

Compactness is the defining attribute of high $\Sigma$ and requires no separate explanation. The remaining three features are monotonic consequences of deeper gravitational potentials: higher $\Sigma$ amplifies all three channels simultaneously (see \S2.4).

The AGN hypothesis faces a specific quantitative difficulty. If the broad emission lines in LRDs originate from AGN broad-line regions, the implied SMBH occupation fraction is extremely high. The majority of spectroscopically confirmed LRDs show broad lines, yet most lack X-ray detections, MIR signatures of hot dust tori, and the time budget for forming abundant SMBHs at $z \sim 7$--$9$ is in tension with $\Lambda$CDM seed formation scenarios (\citealt{Volonteri2010}; \citealt{Latif2016}). The $G_{\rm eff}(\Sigma)$ framework circumvents all of these difficulties: it requires no central engine, only compactness.

\subsection{The Dust Paradox}

The most puzzling aspect of LRDs is not their red color per se, but the \textbf{coexistence of red colors with MIR silence}.

\subsubsection{A Model-Space Artifact}

Current SED fitting codes lack gravitational redshift as a free parameter---they fit spectral shapes using only stellar population age, metallicity, and dust extinction. When a redshift effect from the gravitational potential is not in the model space, the fitting procedure attributes all ``too-red'' signal to the only parameter in the library capable of producing redness: $A_V$.

This means the reported $A_V \approx 1.9$ may not represent pure dust extinction, but rather a \textbf{projected value encoding a mixture of gravitational redshift and dust attenuation} within a restricted model space.

\subsubsection{Cosmic Timeline Constraints}

At $z \sim 7$--$9$ (cosmic age ${\sim}550$--$700$\,Myr), the dominant dust production channel---TP-AGB stars---has not yet had time to operate. Core-collapse supernovae can produce dust, but at low yields ($< 0.1\,M_\odot$ per SN) that are easily destroyed by reverse shocks (\citealt{Gall2011}; \citealt{Matsuura2015}).

\subsubsection{A Falsifiable Prediction: The Balmer Decrement Test}

Dust extinction and gravitational redshift produce \textbf{opposite signatures} in the Balmer decrement. If LRD colors are dominated by gravitational redshift, their Balmer decrements should be systematically consistent with Case~B (${\rm H}\alpha/{\rm H}\beta \approx 2.86$), significantly exceeding the predictions of pure dust models ($< 2.86$). A subset of the RUBIES spectroscopic sample (Sample~B, $N = 36$) has both H$\alpha$ and H$\beta$ measurements; this test constitutes a \textbf{decisive spectroscopic diagnostic}.

\subsection{Magnitude Estimate: Can Gravitational Redshift Explain the Color Excess?}

LRDs are redder than coeval normal galaxies by $\Delta{\rm mag} \sim 0.5$--$2.0$\,mag in $F_{444W} - F_{150W}$. We estimate the gravitational redshift contribution through a four-layer model of increasing physical realism:

\begin{table}[h]
\caption{Four-Layer Gravitational Redshift Model}
\label{tab:four_layer}
\centering
\small
\resizebox{\columnwidth}{!}{\begin{tabular}{llcc}
\hline
Layer & Model & Median $z_{\rm grav}$ & Median $\Delta{\rm mag}$ \\
\hline
1 & Uniform sphere, surface & $6.9 \times 10^{-6}$ & $1.5 \times 10^{-5}$\,mag \\
2 & + emission depth ($\alpha R$) & $1.0 \times 10^{-5}$ & $2.2 \times 10^{-5}$\,mag \\
3 & + S\'ersic $n{=}3$ integration & $3.3 \times 10^{-5}$ & $7.2 \times 10^{-5}$\,mag \\
\textbf{4} & \textbf{+ $G_{\rm eff}(\Sigma)$ ($\varepsilon_g = 1.32$)} & $\mathbf{3.1 \times 10^{-4}}$ & $\mathbf{6.8 \times 10^{-4}}$\,\textbf{mag} \\
\hline
\end{tabular}}
\end{table}

For the 90th-percentile LRD, Layer~4 yields $\Delta{\rm mag} = 3.1 \times 10^{-4}$\,mag, with $G_{\rm eff}/G_N \approx 2.6$ for that source. While the absolute gravitational redshift contribution (${\sim}10^{-4}$\,mag) is far smaller than the observed ${\sim}0.5$\,mag color excess, it is the \textbf{relative ordering}---the monotonic increase of $\Delta{\rm mag}$ with $\Sigma$---that matters for the statistical detection. The four-layer model demonstrates that $G_{\rm eff}(\Sigma)$ introduces a $\Sigma$-dependent color modulation that is absent under standard gravity, and the \S4 statistical tests (Spearman, KS, partial correlation) are sensitive precisely to this monotonic ordering, not to the absolute magnitude. The nonlinear density dependence of $G_{\rm eff}$ is the key enabler: under pure Newtonian gravity (Layers~1--3), the gravitational redshift contribution is negligible and indistinguishable from zero. It is the density-dependent enhancement---independently constrained by the \S4 statistical evidence---that produces a detectable $\Sigma$-color correlation.

The observed color excess is a composite signal: gravitational redshift provides the $\Sigma$-dependent ordering, while other mechanisms (dust, metallicity, stellar population age) contribute the bulk amplitude. This composite picture naturally resolves the MIR silence puzzle: the reported $A_V \approx 1.9$ may encode a gravitational redshift contribution within SED models that lack a non-dust reddening channel, reducing the true dust requirement. Complete calculations, including Monte Carlo error propagation, are presented in Appendix~E.

%% ============================================================
%%  §6  DISCUSSION
%% ============================================================

\section{Discussion}

\subsection{Independent Confirmation: Euclid Compact Candidates}

\citealt{Tumborang2026} identified 233 V-shaped SED sources at $z > 4$ in the 0.75\,deg$^2$ COSMOS field using Euclid (NISP/VIS + Spitzer/IRAC), of which 16 are classified as compact ``amplified LRD'' candidates with stellar masses up to $3 \times 10^{10}\,M_\odot$---approximately an order of magnitude more massive than the typical LRDs in the primary sample (Sample~A).

This Euclid detection supports the $G_{\rm eff}(\Sigma)$ framework along three independent dimensions:

1.~\textbf{Instrumental independence.} Euclid and JWST employ fundamentally different detectors, optical systems, and calibration pipelines.

2.~\textbf{Mass-range extrapolation.} The Euclid candidates extend to higher stellar masses. Under $G_{\rm eff}(\Sigma)$, these more massive systems should exhibit stronger gravitational redshift signatures, with predicted $\Delta{\rm mag} \sim 0.1$--$0.6$\,mag depending on their surface densities.

3.~\textbf{Low AGN contamination.} \citealt{Tumborang2026} estimate an AGN fraction of $< 10\%$ among their compact candidates.

\subsection{Cross-Validation: Signal Survival Under Sample Contamination (JADES DR5)}

\subsubsection{Motivation}

The COSMOS-based analysis (\S4) uses the \citealt{Kokorev2024} photometric catalog, which applies specific color and morphological cuts to define LRDs. We address robustness by applying the framework to the JADES DR5 deep fields (GOODS-North and GOODS-South), using the broader criteria of \citealt{Rinaldi2026}.

\subsubsection{Results}

I performed the same statistical tests applied in \S4.1 (Table~\ref{tab:jades_xval}):

\begin{table*}[!t]
\caption{COSMOS vs.\ JADES Cross-Validation}
\label{tab:jades_xval}
\centering
\footnotesize
\begin{tabular}{lccc}
\hline
Test & COSMOS ($N = 260$) & JADES ($N = 1{,}467$) & Comparison \\
\hline
Raw $\rho$ & $+0.505$ & $+0.505$ & Identical \\
$\rho_p$ (ctrl $z$) & $+0.411$ (6.8$\sigma$) & $-0.172$ (6.6$\sigma$) & \textbf{Sign reversal} \\
$\rho_p$ (ctrl $z + L$) & $\mathbf{+0.341}$ (5.6$\sigma$) & $\mathbf{+0.136}$ (5.2$\sigma$) & Direction consistent \\
KS $D$ (high vs.\ low $\Sigma$) & 0.574 (6.2$\sigma$) & 0.431 (16.6$\sigma$) & Consistent \\
\hline
\end{tabular}
\end{table*}

When both redshift and luminosity are controlled, the JADES signal recovers to $\rho_p = +0.136$ (5.2$\sigma$), consistent in direction with the COSMOS result but attenuated by ${\sim}60\%$.

\begin{figure}[!t]
\centering
\includegraphics[width=\columnwidth]{Fig6_JADES_vs_COSMOS.png}
\caption{JADES DR5 cross-validation. Despite ${\sim}83\%$ sample contamination (82.6\% of candidates are blue), the $\Sigma$--color signal survives in direction ($\rho_p = +0.136$, 5.2$\sigma$). Signal attenuation ($f_{\rm true} \approx 0.40$) quantitatively matches independent purity estimates.}
\label{fig:jades}
\end{figure}

\subsubsection{Contamination Diagnosis}

The JADES candidates reveal severe contamination: only 255/1,467 (17.4\%) satisfy $\log_{10}(F_{444W}/F_{150W}) > 0$, compared to 100\% in the COSMOS sample. This is independently confirmed by \citealt{Rinaldi2026}, who found that only ${\sim}39$--$50\%$ of photometric candidates pass visual inspection.

\subsubsection{Quantitative Consistency of Signal Dilution}

In a contaminated sample, $\rho_{\rm obs} \approx f_{\rm true} \times \rho_{\rm true}$. Substituting $0.136 = f_{\rm true} \times 0.341$ yields $f_{\rm true} \approx 0.40$, in quantitative agreement with the \citealt{Rinaldi2026} purity estimate of 39\%--50\%.

\subsection{Limitations, Falsifiable Predictions, and Falsification Conditions}

\subsubsection{Limitations}

We explicitly acknowledge four principal limitations:

1.~\textbf{Systematic uncertainties in $\Sigma$ measurements.} Different SED fitting codes and PSF deconvolution methods introduce absolute $\Sigma$ offsets of $\pm 0.3$--$0.5$\,dex (\S2.5).

2.~\textbf{Photometric redshift precision.} Sample~A relies on $z_{\rm phot}$ with uncertainties $\Delta z/(1+z) \sim 0.03$--$0.05$.

3.~\textbf{Selection function.} The 260 LRDs are drawn from JWST deep fields, which may not represent the full LRD population.

4.~\textbf{Parameter degeneracy.} The two free parameters ($\varepsilon_g$, $\beta$) are constrained within ${\sim}2$\,dex of $\Sigma$ dynamic range.

\subsubsection{Falsifiable Prediction Matrix}

\begin{table}[h]
\caption{Falsifiable Predictions of the $G_{\rm eff}(\Sigma)$ Framework}
\label{tab:predictions}
\centering
\scriptsize
\setlength{\tabcolsep}{3pt}
\begin{tabular}{@{}l l l@{}}
\hline
\# & Prediction & Status \\
\hline
P1 & H$\alpha$/H$\beta$ $\approx$ Case~B (2.86) & Pending \\
P2 & Euclid compact candidates: positive $\Sigma$--color & Pending \\
P3 & Higher-$\Sigma$ $\to$ broader lines (larger samples) & Pending \\
P4 & Low-$\Sigma$ (CEERS-1025): VMS permitted ($> 300\,M_\odot$) & Pending \\
P5 & High-$\Sigma$ (SPURS): IMF truncation ($\lesssim 90\,M_\odot$) & Pending \\
\hline
\end{tabular}
\end{table}

\subsubsection{Falsification Conditions}

The framework is explicitly falsifiable through four channels:

1.~\textbf{Reversal rate convergence.} If the 34.6\% reversal rate converges to ${\sim}50\%$ in larger samples, the signal is a statistical artifact.

2.~\textbf{Sign reversal.} If $\rho_p(\Sigma,\,{\rm color})$ reverses sign in larger samples, the density-driven hypothesis is falsified.

3.~\textbf{Balmer decrement (decisive binary test).} If the median H$\alpha$/H$\beta$ of confirmed LRDs is significantly below Case~B (${\rm H}\alpha/{\rm H}\beta < 2.5$), gravitational redshift cannot be the \textbf{dominant} contributor. Conversely, if H$\alpha$/H$\beta \approx 2.86$ within observational uncertainty, the dust extinction contribution is limited to $A_V \lesssim 0.3$\,mag.

4.~\textbf{IMF dichotomy (P4/P5 binary test).} The Jeans mass scaling $M_J \propto G_{\rm eff}^{-3/2}$ predicts: high-$\Sigma$ sources should show \textbf{suppressed} VMS, low-$\Sigma$ sources should \textbf{permit} VMS. An inverse pattern would directly falsify the framework.

\subsubsection{Summary: Three Levels of Evidence}

\begin{table}[h]
\caption{Three Levels of Evidence}
\label{tab:evidence_levels}
\centering
\footnotesize
\begin{tabular}{lll}
\hline
Tier & Question & Status \\
\hline
\textbf{I} & Does $\Sigma$ correlate with color? & Established (5.6$\sigma$) \\
\textbf{II} & Is $\Sigma$ an independent factor? & Established (34.6\%, $p < 10^{-30}$) \\
\textbf{III} & Is $G_{\rm eff}(\Sigma)$ the explanation? & Best candidate \\
\hline
\end{tabular}
\end{table}

Tiers~I and II are purely statistical facts. Tier~III is the physical interpretation---$G_{\rm eff}(\Sigma)$ is currently the only framework that simultaneously satisfies the Tier~I+II constraints, but it is not the only possible explanation.

%% ============================================================
%%  §7  CONCLUSION
%% ============================================================

\section{Conclusion}

I have presented the most systematic surface density analysis of JWST Little Red Dots (LRDs) to date, using a sample of 260 photometrically confirmed LRDs and 36 spectroscopically verified counterparts. My principal findings are:

1.~\textbf{Surface density $\Sigma$ is a robust and independent driver of LRD spectral anomalies.} The partial correlation between $F_{444W}/F_{150W}$ and $\log\tilde\Sigma$ is $\rho_p = +0.341$ (5.6$\sigma$), surviving simultaneous control for redshift and bolometric luminosity. Nine independent robustness tests (R1--R9) unanimously preserve the signal, supplemented by formal information-criteria model comparison (Appendix~G) and Monte Carlo systematic error propagation (Appendix~H). The 1,626-pair exhaustive control experiment demonstrates a 34.6\% Type~III reversal rate ($p < 10^{-30}$), decisively falsifying pure Balmer break dominance and establishing $\Sigma$ as an independent causal factor. Spectroscopic verification through the FWHM--$\Sigma$ correlation ($\rho = +0.632$, 4.1$\sigma$) closes an independent evidence chain with four degrees of separation from the photometric analysis.

2.~\textbf{The $G_{\rm eff}(\Sigma)$ framework provides a unified, mechanism-agnostic interpretation.} A single variable---surface density---simultaneously accounts for the four hallmark LRD features: compactness (the definition of high $\Sigma$), red color (gravitational redshift from deepened potential wells), broad emission lines (enhanced virial velocities), and MIR silence (no dust thermal re-emission required). The dust paradox---$A_V \approx 1.9$ coexisting with extreme MIR faintness---dissolves naturally: the reported $A_V$ may encode a gravitational redshift contribution within a restricted SED model space rather than representing pure dust extinction. A four-layer gravitational redshift model (Appendix~E) confirms that $G_{\rm eff} > G_N$ produces a monotonic $\Sigma$-dependent color modulation ($G_{\rm eff}/G_N = 2.0$--$3.6$ across the sample), providing the physical mechanism for the detected $\Sigma$--color ordering.

3.~\textbf{Independent verification strengthens the result.} Euclid-detected compact candidates (\citealt{Tumborang2026}) confirm the $\Sigma$-dependent population with an independent instrument and selection function. The JADES DR5 cross-validation recovers a positive $\Sigma$--color signal ($\rho_p = +0.136$, 5.2$\sigma$) even under ${\sim}83\%$ sample contamination, with signal attenuation quantitatively matching independently estimated purity fractions.

I emphasize that the parameter constraints on $G_{\rm eff}(\Sigma)$ are derived exclusively from the high-$z$ LRD population. Extrapolation to the local universe or other astrophysical populations requires independent observational calibration at each redshift (\S6.3.1).

Five falsifiable predictions (P1--P5) are offered for future testing. Verification of P1, P4, and P5 is within reach of JWST Cycle~4+ programs, Euclid NIRSpec, and infrared spectrographs on the coming generation of 30\,m-class telescopes; their outcomes will directly confirm or falsify this framework. Regardless of whether $G_{\rm eff}(\Sigma)$'s underlying physical mechanism is ultimately validated, this paper establishes a more fundamental conclusion:

\textbf{Surface density $\Sigma$ is the organizing principle of LRD spectral properties.} Any framework claiming to fully explain LRDs must incorporate $\Sigma$ as an independent degree of freedom. The $G_{\rm eff}(\Sigma)$ framework is currently the only interpretation anchored to this observational fact from the outset.

\acknowledgments

I thank the Kokorev et al.\ (2024) team for making the DJA LRD photometric catalog publicly available, and the RUBIES collaboration for providing NIRSpec spectroscopic redshifts and emission-line measurements. I acknowledge the JADES and CEERS survey teams for their open-access JWST data products, which enabled the cross-validation analyses presented in this work. I am grateful to the anonymous reviewer of the earlier ApJL submission (AAS76108) for constructive feedback that motivated the expansion into this full-length article. As a sole investigator without institutional computing resources, I relied on personal hardware; I thank the open-source community for \texttt{Python}, \texttt{astropy}, \texttt{SciPy}, and \texttt{JWST} calibration pipelines.

\dataavailability{All data used in this analysis are drawn from publicly released JWST survey products. The 260-LRD photometric catalog (Sample~A) and 36-LRD spectroscopic catalog (Sample~B) are provided as machine-readable tables on Zenodo (doi:10.5281/zenodo.19587085). The original Kokorev et al.\ (2024) catalog is publicly available via the DJA. JADES DR5 data products are available via the JWST archive at \url{https://archive.stsci.edu/}. All analysis scripts are available upon reasonable request.}

%% ============================================================
%%  APPENDICES
%% ============================================================

\appendix

%% ---------- Appendix A ----------
\section{Sample Catalog}

The full 260-LRD catalog (Sample~A) and 36-LRD spectroscopic catalog (Sample~B) are provided as machine-readable tables, available at Zenodo (DOI pending) as \texttt{lrd\_260\_full.csv} and \texttt{lrd\_36\_spectroscopic.csv}.

\subsection{Column Definitions}

\textbf{Sample~A columns:} ID (Kokorev catalog), FIELD (JWST pointing), $z_{\rm phot}$, $M_{F444W}$ (AB mag), color ($\log_{10}(F_{444W}/F_{150W})$), $\log\tilde\Sigma$ (observational surface-density proxy; Eq.~\ref{eq:Sigma_proxy}), $r_{\rm eff}$ (galfit). Additional columns include RA, Dec, full six-band NIRCam fluxes, $A_V$, $L_{\rm bol}$, SED $\Delta\chi^2$, and FWHM where measured.

\textbf{Sample~B columns:} srcid (RUBIES), FIELD, $z_{\rm spec}$, FWHM, $M_{F444W}$, $\log\tilde\Sigma$, $r_{\rm eff}$.

\subsection{Statistical Summary}

\begin{table}[h]
\caption{Sample Statistics (Median Split by $\Sigma$)}
\label{tab:median_split}
\centering
\footnotesize
\begin{tabular}{lccc}
\hline
Parameter & Full ($N = 260$) & High-$\Sigma$ & Low-$\Sigma$ \\
\hline
$\log\,\tilde\Sigma$ & $-7.42 \pm 0.98$ & $-6.68$ & $-8.15$ \\
$\log_{10}(F_{444W}/F_{150W})$ & $+0.51 \pm 0.47$ & $+0.76$ & $+0.25$ \\
$z_{\rm phot}$ & $6.42 \pm 1.52$ & $6.18$ & $6.65$ \\
$r_{\rm eff}$ & $121 \pm 63$\,pc & $89$\,pc & $158$\,pc \\
\hline
\end{tabular}
\end{table}

%% ---------- Appendix B ----------
\section{1,626 Matched-Pair Control Experiment}

\subsection{Pair Construction Algorithm}

Pairs were constructed from the combined Sample~A + Sample~B catalog ($N = 296$). The algorithm evaluated all $\binom{296}{2} = 43{,}660$ combinations against four hard thresholds: (i)~same JWST field, (ii)~angular separation $\leq 5'$, (iii)~$|\Delta z| \geq 0.3$, (iv)~$|\Delta\log\tilde\Sigma| \geq 0.2$.

\subsection{Dose--Response Relationship}

Among the 562 Type~III reversed pairs:
\begin{equation}
\rho(|\Delta\log\tilde\Sigma|,\, |\Delta{\rm color}|) = +0.236, \quad p = 4.8 \times 10^{-22}
\end{equation}
Each 1\,dex increase in $|\Delta\log\tilde\Sigma|$ produces ${\sim}0.3$\,mag additional $|\Delta{\rm color}|$ reversal.

%% ---------- Appendix C ----------
\section{Spectroscopic Cross-Validation and Control Samples}

\subsection{RUBIES $\times$ DJA Cross-Matching}

Of 36 RUBIES LRDs, 34 were successfully matched to the DJA galfit catalog (2 lacked DJA coverage). The median RUBIES--DJA $m_{\rm F444W}$ offset is $0.02 \pm 0.15$\,mag, confirming photometric consistency.

\subsection{FWHM--$\Sigma$ Regression}

\begin{equation}
{\rm FWHM} = (590 \pm 130) + (311 \pm 75)\,\log\tilde\Sigma \quad {\rm km\,s^{-1}}
\end{equation}
with $R^2 = 0.40$. A 1\,dex increase in $\Sigma$ produces ${\sim}310$\,km\,s$^{-1}$ increase in line width.

\subsection{Non-LRD Control Group}

260 matched control galaxies (non-LRD) have null color--$\Sigma$ correlation: $\rho = +0.04$ ($p = 0.52$, $< 1\sigma$), confirming the signal is specific to the high-$\Sigma$ LRD population.

%% ---------- Appendix D ----------
\section{Comparison with Alternative LRD Models}

\subsection{Model Landscape}

\begin{table*}
\caption{Survey of Proposed LRD Physical Models}
\label{tab:models}
\centering
\scriptsize
\begin{tabular}{clll}
\toprule
\# & Model & Rep.\ Works & $\Sigma$-Dep.\ Color? \\
\midrule
1 & Super-Eddington slim disk & Pacucci+24 & No \\
2 & Gas slab obscuration (AGN) & Inayoshi+24 & No \\
3 & BH-envelope (BH*) & Liu+25,26 & No \\
4 & AGN + stellar composite & JADES teams & No \\
5 & Metal-poor massive seeds & Ivey+26 & No \\
6 & Extreme value statistics & Heather+26 & No \\
7 & SMS pulsation shell & Nandal+26 & No \\
8 & Runaway collision EMS & Rantala+26 & Marginal$^{\dagger}$ \\
9 & Paschen jump / nebular & Sneppen+26 & No \\
10 & PBH seeds & Zhang+26 & No \\
11 & Selection effect & Rinaldi+26 & No \\
12 & Brown dwarf contamination & arXiv:2604.23668 & No \\
\bottomrule
\multicolumn{4}{l}{$^{\dagger}$\scriptsize Runaway collision naturally favors high-$\Sigma$ environments but no quantitative $\Sigma$--color prediction exists in the published literature.} \\
\end{tabular}
\end{table*}

\textbf{Key observation}: Models~1--7 and 9--12 explain \textit{why LRDs exist} but none predicts \textit{why more compact LRDs are systematically redder}---the structural question addressed by the $5.6\sigma$ correlation.

\subsection{Cross-Model Falsification Implications}

The Balmer decrement test (P1) constrains models~1--4 to a subdominant role. The VMS dichotomy test (P4/P5) creates tension for model~8 (runaway collision), which requires high density to trigger but would predict VMS precisely where $G_{\rm eff}(\Sigma)$ predicts IMF truncation.

%% ---------- Appendix E ----------
\section{Gravitational Redshift Magnitude Estimate}

\subsection{Physical Framework}

The gravitational redshift for a photon escaping from radius $r$ within a spherically symmetric mass distribution is:
\begin{equation}
z_{\rm grav} = \frac{\Delta\phi}{c^2} = \frac{1}{c^2}\left[\phi(r_{\rm emit}) - \phi(r_{\rm obs})\right]
\end{equation}
For $z_{\rm grav} \ll 1$, the magnitude shift at a single wavelength is $\Delta m \approx 1.086\,z_{\rm grav}$. For a broadband color index (e.g., $F_{444W} - F_{150W}$), the shift in each filter band adds constructively because gravitational redshift dims the longer-wavelength band (where the source emits near the spectral peak) while leaving the shorter-wavelength band (dominated by the Rayleigh--Jeans tail) relatively unaffected. The resulting color excess is therefore $\Delta{\rm color} = \Delta m_{444} - \Delta m_{150} \approx 2\!\times\!1.086\,z_{\rm grav} = 2.172\,z_{\rm grav}$.

\subsection{Four-Layer Model}

\textbf{Layer 1} (Uniform sphere, surface emission): $z_1 = G_N M/(R c^2)$; for a representative LRD ($M_* = 10^{10}\,M_\odot$, $R = 100$\,pc, $\Sigma = 3.2 \times 10^5\;M_\odot\,{\rm pc}^{-2}$): $z_1 = 4.8 \times 10^{-6}$, $\Delta{\rm mag}_1 = 1.0 \times 10^{-5}$\,mag.

\textbf{Layer 2} (+ emission depth $\alpha R$, $\alpha = 0.7$): $z_2 = 6.0 \times 10^{-6}$, $\Delta{\rm mag}_2 = 1.3 \times 10^{-5}$\,mag ($\times 1.26$ enhancement over surface emission).

\textbf{Layer 3} (+ S\'ersic $n = 3$, emission at $\alpha = 0.2\,R_e$, line-of-sight integration, geometric factor $\times 1.3$): $z_3 = 2.3 \times 10^{-5}$, $\Delta{\rm mag}_3 = 5.0 \times 10^{-5}$\,mag ($\times 4.8$ cumulative enhancement).

\textbf{Layer 4} (+ $G_{\rm eff}(\Sigma)$ with fitted parameters $\varepsilon_g = 1.32$, $\beta = 0.097$, $\Sigma_0 = 1.6 \times 10^5\;M_\odot\,{\rm pc}^{-2}$; Table~\ref{tab:geff_params}): $G_{\rm eff}/G_N = 2.41$ for the representative source ($\Sigma/\Sigma_0 = 2.0$); $z_4 = 5.5 \times 10^{-5}$, $\Delta{\rm mag}_4 = 1.2 \times 10^{-4}$\,mag ($\times 11.5$ cumulative enhancement; $\times 2.4$ from $G_{\rm eff}$ alone).

\begin{figure}[!t]
\centering
\includegraphics[width=\columnwidth]{FigE1_Magnitude_DeepDive.png}
\caption{Four-layer gravitational redshift magnitude estimate for a representative LRD ($M_* = 10^{10}\,M_\odot$, $R = 100$\,pc). Each layer adds physical realism. The $G_{\rm eff}(\Sigma)$ enhancement (Layer~4) increases the gravitational redshift by a factor of ${\sim}2.4{\times}$ relative to standard gravity (Layer~3). See Table~\ref{tab:grav_z} for the full 260-LRD sample statistics.}
\label{fig:magnitude}
\end{figure}

\subsection{Full Sample Results}

\begin{table}[h]
\caption{Gravitational Redshift Distribution Across 260 LRDs}
\label{tab:grav_z}
\centering
\footnotesize
\begin{tabular}{lccc}
\hline
Statistic & Layer~1 & Layer~3 & \textbf{Layer~4} \\
\hline
Median $\Delta{\rm mag}$ & $4.9 \times 10^{-6}$ & $2.3 \times 10^{-5}$ & $5.4 \times 10^{-5}$ \\
90th percentile & $2.5 \times 10^{-5}$ & $1.2 \times 10^{-4}$ & $3.1 \times 10^{-4}$ \\
95th percentile & $3.8 \times 10^{-5}$ & $1.8 \times 10^{-4}$ & $4.9 \times 10^{-4}$ \\
\hline
\end{tabular}
\end{table}

\subsection{Monte Carlo Error Propagation}

Monte Carlo error propagation (1,000 iterations, $M_*$ perturbed by $\pm 0.2$\,dex, $r_{\rm eff}$ by $\pm 30\%$) confirms that the Layer~4 values are robust: the 90th-percentile $\Delta{\rm mag}$ remains in the range $(2.1$--$4.0) \times 10^{-4}$\,mag in 95\% of realizations, with $G_{\rm eff}/G_N$ varying between 2.0 and 3.6 across the $G_{\rm eff}(\Sigma)$ parameter uncertainty ($\varepsilon_g = 1.32 \pm 0.08$, $\beta = 0.097 \pm 0.030$).

%% ---------- Appendix F ----------
\section{SPURS/CEERS IMF Analysis}

\subsection{The Jeans Mass Argument Under Enhanced Gravity}

Under $G_{\rm eff}(\Sigma) > G_N$, the Jeans mass scales as $M_J \propto G_{\rm eff}^{-3/2}$: enhanced gravity lowers the fragmentation mass threshold, suppressing very massive star (VMS) formation.

\begin{table}[h]
\caption{IMF Upper Limits for Five High-$z$ Compact Sources}
\label{tab:imf_limits}
\centering
\footnotesize
\begin{tabular}{lcccc}
\hline
Target & $z$ & $\log\Sigma$ & $G_{\rm eff}/G_N$ & IMF cap ($M_\odot$) \\
\hline
SPURS-1 & $9.31$ & 4.90 & 2.23 & $\sim 90$ \\
CEERS-1019-C & 8.69 & 4.43 & 2.11 & $\sim 98$ \\
CEERS-1019-A & 8.69 & 4.07 & 2.02 & $\sim 104$ \\
CEERS-1019-B & 8.69 & 3.94 & 1.99 & $\sim 107$ \\
CEERS-1025 & 8.71 & 2.42 & 1.71 & $\sim 134$ \\
\hline
\end{tabular}
\end{table}

The $G_{\rm eff}(\Sigma)$ parameters ($\varepsilon_g = 1.317$, $\beta = 0.097$, $\Sigma_0 = 1.6 \times 10^5\;M_\odot\;{\rm pc}^{-2}$; Table~\ref{tab:geff_params} reports $\varepsilon_g$ rounded to 1.32) are those fitted to the 260-LRD sample in \S4. The IMF upper limits are derived under the assumption of a VMS-permitting baseline of $M_{\rm up} = 300\,M_\odot$ in the Newtonian limit ($G_{\rm eff} = G_N$); the tabulated caps represent the reduced $M_J$ under the locally enhanced $G_{\rm eff}$.

\begin{figure}[!t]
\centering
\includegraphics[width=\columnwidth]{FigF1_SPURS_IMF.png}
\caption{IMF upper limits vs.\ surface density for five high-$z$ compact sources. Higher $\Sigma$ produces lower IMF caps, predicting systematic VMS suppression at high density.}
\label{fig:imf}
\end{figure}

%% ---------- Appendix G ----------
\section{Model Comparison via Information Criteria}

\subsection{Methodology}

I compute AIC and BIC for three competing models on the filtered color--$\Sigma$ dataset ($n = 253$; see \S3):
\begin{equation}
{\rm AIC} = n\ln\!\left(\frac{{\rm RSS}}{n}\right) + 2k, \qquad {\rm BIC} = n\ln\!\left(\frac{{\rm RSS}}{n}\right) + k\ln n
\end{equation}

\subsection{Competing Models}

\textbf{M0 (null)}: ${\rm color} = c_0$ ($k = 1$). \textbf{M1 (dust/AGN)}: ${\rm color} = a_0 + a_1 z_{\rm phot} + a_2 \log L_{\rm bol}$ ($k = 3$). \textbf{M2 [$G_{\rm eff}(\Sigma)$]}: ${\rm color} = b_0 + b_1 \log\tilde\Sigma + b_2 z_{\rm phot} + b_3 \log L_{\rm bol}$ ($k = 4$).

\subsection{Results}

\begin{table}[h]
\caption{Information Criteria and Cross-Validation}
\label{tab:info_criteria}
\centering
\footnotesize
\begin{tabular}{lccccccc}
\hline
Model & $k$ & RSS & $\Delta$AIC & $\Delta$BIC & $R^2$ & $Q^2$ & Drop \\
\hline
M0 (null) & 1 & 19.0 & +49.1 & +38.5 & 0.000 & $-$0.008 & --- \\
M1 (dust/AGN) & 3 & 16.5 & +18.2 & +14.6 & 0.129 & 0.103 & 20.5\% \\
\textbf{M2 ($G_{\rm eff}$)} & \textbf{4} & \textbf{15.3} & \textbf{0} & \textbf{0} & \textbf{0.196} & \textbf{0.163} & \textbf{16.8\%} \\
\hline
\end{tabular}
\end{table}

$G_{\rm eff}(\Sigma)$ achieves $\Delta{\rm AIC} = 49.1$ and $\Delta{\rm BIC} = 38.5$ vs.\ null, and $\Delta{\rm AIC} = 18.2$, $\Delta{\rm BIC} = 14.6$ vs.\ standard dust/AGN. The additional $\Sigma$ parameter reduces RSS by 7.7\% relative to M1 and 19.6\% relative to M0. BIC's heavier parameter penalty still favors M2, confirming the improvement is not driven by overfitting.

LOO-CV: $Q^2 = 0.163$ for M2 drops 16.8\% from in-sample $R^2$, compared to 20.5\% for M1. M2 maintains positive predictive power while consistently outperforming M1, confirming $\Sigma$ provides genuinely independent predictive information.

%% ---------- Appendix H ----------
\section{Systematic Error Budget}

The central result $\rho_p = +0.341$ (5.6$\sigma$) must withstand scrutiny of measurement systematics. I quantify three dominant error sources using Monte Carlo perturbation (1,000 iterations per test).

\subsection{Per-Source Perturbation Results}

\begin{table}[h]
\caption{Impact of Per-Source Systematic Perturbations on $\rho_p$}
\label{tab:mc_perturbation}
\centering
\footnotesize
\begin{tabular}{lcccc}
\hline
Source & $\sigma$ & Median $\rho_p$ & 95\% CI & Key Finding \\
\hline
$r_{\rm eff}$ & 30\% & +0.226 & [+0.132, +0.316] & Largest risk \\
$z_{\rm phot}$ & 0.10 & +0.342 & [+0.336, +0.346] & No effect \\
SED ($L_{\rm bol}$) & 0.30\,dex & +0.357 & [+0.334, +0.378] & Amplifies \\
\hline
\end{tabular}
\end{table}

$r_{\rm eff}$ noise is the dominant systematic. Counter-intuitively, SED code noise slightly amplifies the signal: added $L_{\rm bol}$ noise decorrelates the control variable, reducing confounding absorption.

\subsection{Combined Worst Case}

Under the conservative combined scenario ($r_{\rm eff}$ 30\%, $z_{\rm phot}$ 0.10, $L_{\rm bol}$ 0.30\,dex): median $\rho_p = +0.239$ (3.8$\sigma$); 97.5\% of realizations exceed 2.4$\sigma$.

\textbf{Conclusion}: The 5.6$\sigma$ signal is robust to all realistic systematic perturbations. The dominant vulnerability is $r_{\rm eff}$ measurement noise, which attenuates but does not reverse the signal. Improving NIRCam PSF models will strengthen the detection.



\begin{thebibliography}{99}

\bibitem[Adams et al.(2024)]{Adams2024}
Adams, N.\ J., et al.\ 2024, ApJ, 965, 169

\bibitem[Carnall et al.(2018)]{Carnall2018}
Carnall, A.\ C., et al.\ 2018, MNRAS, 480, 4379

\bibitem[Carnall et al.(2024)]{Carnall2024}
Carnall, A.\ C., et al.\ 2024, MNRAS, 534, 325

\bibitem[Chevallier \& Polarski(2001)]{Chevallier2001}
Chevallier, M., \& Polarski, D.\ 2001, Int.\ J.\ Mod.\ Phys.\ D, 10, 213

\bibitem[Fujimoto et al.(2023)]{Fujimoto2023}
Fujimoto, S., et al.\ 2023, ApJL, 949, L25

\bibitem[Gall et al.(2011)]{Gall2011}
Gall, C., et al.\ 2011, A\&A, 528, A13

\bibitem[Hviding et al.(2024)]{Hviding2024}
Hviding, R.\ E., et al.\ 2024, AJ, 168, 220

\bibitem[Kocevski et al.(2025)]{Kocevski2024}
Kocevski, D.\ D., et al.\ 2025, ApJ, 986, 126

\bibitem[Kokorev et al.(2024)]{Kokorev2024}
Kokorev, V., et al.\ 2024, ApJ, 968, 38

\bibitem[Labbe et al.(2023)]{Labbe2023}
Labbe, I., et al.\ 2023, Nature, 616, 266

\bibitem[Latif \& Ferrara(2016)]{Latif2016}
Latif, M.\ A., \& Ferrara, A.\ 2016, PASA, 33, e051

\bibitem[Leja et al.(2019)]{Leja2019}
Leja, J., et al.\ 2019, ApJ, 876, 3

\bibitem[Linder(2005)]{Linder2005}
Linder, E.\ V.\ 2005, Phys.\ Rev.\ D, 72, 043529

\bibitem[Matthee et al.(2024)]{Matthee2024}
Matthee, J., et al.\ 2024, ApJ, 963, 129

\bibitem[Matsuura et al.(2015)]{Matsuura2015}
Matsuura, M., et al.\ 2015, ApJ, 800, 50

\bibitem[Peng et al.(2010)]{Peng2010}
Peng, C.\ Y., et al.\ 2010, AJ, 139, 2097

\bibitem[Perez-Gonzalez et al.(2023)]{PerezGonzalez2023}
Perez-Gonzalez, P.\ G., et al.\ 2023, ApJL, 951, L1

\bibitem[Rinaldi et al.(2026)]{Rinaldi2026}
Rinaldi, P., et al.\ 2026, arXiv:2604.07138

\bibitem[Tumborang et al.(2026)]{Tumborang2026}
Tumborang, C.\ F., et al.\ 2026, arXiv:2604.16178

\bibitem[Uehling(1935)]{Uehling1935}
Uehling, E.\ A.\ 1935, Phys.\ Rev., 48, 55

\bibitem[Venturi et al.(2024)]{Venturi2024}
Venturi, G., Carniani, S., Parlanti, E., et al.\ 2024, A\&A, 691, A19

\bibitem[Volonteri(2021)]{Volonteri2021}
Volonteri, M.\ 2021, Nat.\ Rev.\ Phys., 3, 732

\bibitem[Volonteri(2010)]{Volonteri2010}
Volonteri, M.\ 2010, A\&ARv, 18, 279

\end{thebibliography}

\end{document}
